\section{Main results}
\label{sec:main-results}

\subsection{Patch representation}

For~$A\Subset\Lambda$, Theorems~\ref{thm:patch-factorization} and
\ref{thm:patch-representation} reorganize the signed Feynman--Kac
representation of~$P_t\chi_A$. Conditional on the successful-interaction
skeleton, the randomness in distinct patches is independent. Averaging it
patch by patch gives an exact representation split between bulk patches
completed before time~$t$ and end patches reaching time~$t$. The bulk product
records the history, while the end product also depends on the terminal
configuration. Cancellations are therefore taken before any absolute-value
estimate is made.

Patch positivity means that every bulk-patch contribution obtained from this
averaging is nonnegative. It has a local coefficient characterization. If~$c_i(\mb 0)+c_i(\mb{0}^{i})>0$,
Theorem~\ref{thm:patch-positivity-criterion} says that patch positivity is
equivalent to
\begin{equation*}
c_i^0(S)+c_i^1(S)\leq0,
\qquad
c_i^1(\vn)c_i^0(S)\geq c_i^0(\vn)c_i^1(S)
\qquad
\text{for every }\vn\neq S\subseteq N(i);
\end{equation*}
if~$c_i(\mb 0)+c_i(\mb{0}^{i})=0$, it is equivalent to~$c_i\equiv0$. Thus patch positivity can be
checked directly from the multilinear coefficients of the two directional
flip rates.

End patches retain the terminal state and require a one-density threshold.
Proposition~\ref{prop:critical-profile-formula} identifies the minimal
profile~$\mb p^\star$ for which all end-patch contributions are nonnegative:
\begin{equation*}
p_i^\star
=
\max\cb{
0,
\sup_{\substack{
\vn\neq S\subseteq N(i)\\
c_i^0(S)+c_i^1(S)<0
}}
\frac{c_i^0(S)}{c_i^0(S)+c_i^1(S)}
}.
\end{equation*}

\subsection{Centered-moment order}

The profile~$\mb p^\star$ determines the centered moments on which the
calm--facilitating order survives. Set
$\chi_A^*(\eta)=\prod_{i\in A}(\eta(i)-p_i^\star)$ and
write~$\mu\preceq_*\nu$ when~$\mu(\chi_A^*)\leq\nu(\chi_A^*)$ for every
$A\Subset\Lambda$. Theorem~\ref{thm:centered-moment-order} gives
\begin{equation*}
\mu\preceq_*\nu
\quad\Longrightarrow\quad
\mu P_t\preceq_*\nu P_t
\qquad
\text{for every }t\geq0,
\end{equation*}
without requiring either measure to have nonnegative centered moments. In
particular, the class
\begin{equation*}
\mathcal M_*
=
\cb{
\mu:\mu\bb{\chi_A^*}\geq0
\text{ for every }A\Subset\Lambda
}
\end{equation*}
is preserved. The dynamics need not preserve an order of configurations, but
it preserves the order of all products of deviations from the critical
calm-state profile.

The patch representation also compares the system with one obtained by
removing environment-independent~$1\to0$ transitions. If~$P_t'$ denotes the
semigroup after removing such pure deaths, then
Theorem~\ref{thm:pure-death-comparison} gives
\begin{equation*}
(\mu P_t)(\chi_A)
\leq
(\mu P_t')(\chi_A)
\qquad
\text{for every }\mu\in\mathcal M_*.
\end{equation*}

\subsection{Common invariant limit}

Assume that~$\Lambda$ has polynomial growth of exponent~$D$ and that the spin
system contains pure-death noise at rate~$\varepsilon>0$. For~$K\geq0$, let
\begin{equation*}
\mathcal M_{-,K}
=
\cb{
\mu:
\frac{\mu+K\mu_{\mb 1}}{1+K}\in\mathcal M_*
},
\qquad
\mathcal M_-=\bigcup_{K\geq0}\mathcal M_{-,K}.
\end{equation*}
This class contains Bernoulli product measures with sufficiently high
calm-state density and point masses on configurations with finitely many
facilitating states. Theorem~\ref{thm:common-invariant-limit} gives an invariant
measure~$\pi$ such that, for every local function~$f$,
\begin{equation*}
\sup_{\mu\in\mathcal M_{-,K}}
\abs{(\mu P_t)(f)-\pi(f)}
\leq
(1+2K)K_f(1+t)^D e^{-\varepsilon t/2}
\end{equation*}
for every~$K\geq0$ and~$t\geq0$. The limit is obtained without identifying an
invariant measure in advance. In the notation of
Section~\ref{sec:representation}, its monomial moments are
\begin{equation*}
\pi(\chi_A)
=
\E_A\Cb{
\prod_{P\in\mathcal P}C(P)
\ind\bb{\abs{\mathcal P}<\infty}
}.
\end{equation*}
If~$\mb p^\star\leq\mbs{\tfrac12}$, then every probability measure belongs
to~$\mathcal M_{-,1}$,~$\pi$ is the unique invariant measure, and
Corollary~\ref{cor:uniform-ergodicity} gives
\begin{equation*}
\sup_{\eta\in\Omega}
\abs{P_tf(\eta)-\pi(f)}
\leq
K_f(1+t)^D e^{-\varepsilon t/2}.
\end{equation*}

Section~\ref{sec:applications} introduces the contact process, FA-1f, and BABP
and states the consequences of these results for the latter two models. The
signed dual and graphical construction are developed in
Section~\ref{sec:signed-dual}; Sections~\ref{sec:patches} and
\ref{sec:representation} establish patch factorization and the representation
formula. Sections~\ref{sec:patch-positivity}, \ref{sec:moment-order}, and
\ref{sec:convergence} prove the coefficient criterion, the centered-moment
comparisons, and the common invariant-limit theorem.
